\documentclass[11pt]{article}
\usepackage{explorer}
\newcommand{\ExplorerLanguage}{EN}
\newcommand{\ExplorerNumber}{07}
\newcommand{\ExplorerKind}{Lesson and method}
\begin{document}
\ExplorerTitle{High school · probability}{Reversing a conditional probability}
\ExplorerSection{Before you begin}
Probability trees, intersections and conditional probability.\par
\ExplorerSection{Ideas and method}
\(P(A|B)\) is the probability of A given that B occurred. If \(P(B) > 0, P(A|B) = \frac{P(A \cap  B)}{P(B)}\). The information B restricts the sample space.\par
\(P(A \cap  B) = P(A)P(B|A)\). If A and its complement form the first two branches of a tree, \(P(B) = P(A)P(B|A) + P(\mathrm{not} A)P(B|\mathrm{not} A)\).\par
Bayes’ formula follows from these two identities. \(P(A|B)\) generally differs from \(P(B|A)\): the denominators refer to different populations.\par
\ExplorerSection{Worked example}
Choose urn A with probability \(\frac{1}{3}\) and B with probability \(\frac{2}{3}\). A contains 3 red and 1 blue ball; B contains 1 red and 3 blue balls. \(P(\mathrm{red}) = (\frac{1}{3})(\frac{3}{4}) + (\frac{2}{3})(\frac{1}{4}) = \frac{5}{12}\). Given a red ball, \(P(A|\mathrm{red}) = \frac{\frac{1}{4}}{\frac{5}{12}} = \frac{3}{5}\).\par
\begin{center}\begin{tikzpicture}[x=1.6cm,y=.8cm]\draw[mint,thick](0,0)--(1,1)node[midway,above]{$\frac13$};\draw[mint,thick](0,0)--(1,-1)node[midway,below]{$\frac23$};\node at(1.2,1){$A$};\node at(1.2,-1){$B$};\draw[violet,thick](1.5,1)--(3,1.5)node[midway,above]{$\frac34$};\draw[violet,thick](1.5,-1)--(3,-.5)node[midway,above]{$\frac14$};\node[right]at(3,1.5){$A\cap R:\ \frac14$};\node[right]at(3,-.5){$B\cap R:\ \frac16$};\end{tikzpicture}\end{center}
Only branches leading to a red ball are shown.\par
\ExplorerSection{Explore}
Represent the example with 120 hypothetical trials matching the proportions exactly. Count red balls from each urn.\par
\end{document}
